\subsection{Conjugate extensions. Conjugate elements}

DEFINITION 1. --- Let E and F be two extensions of K contained in Ω. We shall say that E and F are conjugate (in Ω) if there exists a K-automorphism u of Ω such that \(u(E) = F\) . Two elements x and y of Ω are said to be conjugate over K if there exists a K-automorphism u of Ω such that \(u(x) = y\) .

Let E and F be two extensions of K contained in \(\Omega\) . By Prop. 1, E and F are conjugate over K if and only if they are isomorphic extensions of K. This is so in particular when there exist two subsets A and B of \(\Omega\) such that \(\mathrm{E} = \mathrm{K}(\mathrm{A})\) and \(\mathrm{F} = \mathrm{K}(\mathrm{B})\) and a K-automorphism u of \(\Omega\) such that \(u(\mathrm{A}) = \mathrm{B}\) .

The relation « x and y are conjugate over K » is an equivalence relation in \(\Omega\) ; the equivalence classes are called the conjugacy classes in \(\Omega\) ; they are the orbits in \(\Omega\) of the group of all K-automorphisms of \(\Omega\) .

PROPOSITION 2. --- Let \(x\) and \(y\) be two elements of \(\Omega\) . Then the following conditions are equivalent :

\begin{enumerate}
\def\labelenumi{\alph{enumi})}
\item
  \(x\) and \(y\) are conjugate over \(\mathbf{K}\) .
\item
  There exists a K-isomorphism \(v\) of \(\mathbf{K}(x)\) onto \(\mathbf{K}(y)\) such that \(v(x) = y\) .
\item
  \(x\) and \(y\) have the same minimal polynomial over \(\mathbf{K}\) .
\end{enumerate}

Suppose first that x and y are conjugate over K; let u be a K-automorphism of \(\Omega\) such that \(u(x) = y\) and let f be the minimal polynomial of x over K. We have

\[
f (y) = f (u (x)) = u (f (x)) = 0,
\]

and \(f\) is a monic irreducible polynomial in \(\mathbf{K}[\mathbf{X}]\) ; therefore (V, p.~16, Th. 1, c)), \(f\) is the minimal polynomial of \(y\) over \(\mathbf{K}\) . Thus \(a)\) implies \(c)\) .

Suppose now that x and y have the same minimal polynomial f over K. There exists a K-isomorphism of the field \(\mathrm{K}[\mathrm{X}]/(f)\) onto \(\mathrm{K}(x)\) (resp. onto \(\mathrm{K}(y)\) ) mapping the residue class of X modulo (f) to x (resp. y) (V, p.~16, Th. 1, b)); hence there exists a K-isomorphism v of \(\mathrm{K}(x)\) onto \(\mathrm{K}(y)\) such that \(v(x)=y\) . Therefore c) implies b).

Finally, under the hypothesis \(b\) ), Prop. 1 implies the existence of a K-automorphism \(u\) of \(\Omega\) extending \(v\) ; we then have \(u(x) = y\) , hence \(x\) and \(y\) are conjugate over \(K\) , so \(b\) ) implies \(a\) ).

COROLLARY 1. --- Let x be an element of Ω of degree n over K, and let f be the minimal polynomial of x over K. The conjugates of x over K are the roots of f in Ω and their number is at most equal to n.

COROLLARY 2. --- Let x be an element of Ω of degree n over K. For x to be separable over K it is necessary and sufficient that x should have n conjugates in Ω; when this is so, all the conjugates of x over K are separable over K.

Let \(f\) be the minimal polynomial of \(x\) over \(K\) ; its roots are the conjugates of \(x\) over \(K\) , and each of these roots admits \(f\) as minimal polynomial over \(K\) . Now \(x\) is separable over \(K\) if and only if the polynomial \(f\) is separable (V, p.~39, Prop. 5), that is, if \(f\) has \(n\) distinct roots in \(\Omega\) . Cor. 2 follows from this.

COROLLARY 3. --- Let G be the group of K-automorphisms of \(\Omega\) . The set of invariants of G in \(\Omega\) is the relative \(p\) -radical closure of K in \(\Omega\) (V, p.~25).

In other words, an element \(x\) of \(\Omega\) is \(p\) -radical over \(\mathbf{K}\) if and only if it has no conjugate other than itself.

Let \(p\) be the characteristic exponent and let \(x \in \Omega\) . By V, p.~44, Prop. 13, there exists an integer \(m \geqslant 0\) such that \(y = x^{p^m}\) is algebraic and separable over K. Now \(x\) is invariant under G if and only if \(y\) is invariant under G; by Cor. 2 this is the case if and only if \(y\) is of degree 1 over K, which amounts to saying that \(y \in K\) . Now the corollary is an immediate consequence.

